\documentclass[12pt]{article}
\usepackage[utf8]{inputenc}
\usepackage[T1]{fontenc}
\usepackage{hyperref}
\title{Vehes Language Resources for K-12 Curriculum Development}
\author{Compiled by Claude Sonnet 5 (claude-sonnet-5) for Andr\'as Kornai}
\date{}
\begin{document}
\maketitle
\section{Introduction}
Vehes (ISO 639-3 val, Glottocode vehe1237) is an endangered Oceanic language of Morobe Province, Papua New Guinea, with about 70 speakers in one village near the coast between Lae and Salamaua. It is also called Buasi or Vehees, and belongs to the Huon Gulf branch. Very little exists.
\section{Primary Educational Resources}
None found.
\section{Linguistic and Reference Materials}
\begin{itemize}
\item Vehes language, Wikipedia: \url{https://en.wikipedia.org/wiki/Vehes_language}
\item Numeral system data (Gillian Sankoff, 1998): \url{https://lingweb.eva.mpg.de/channumerals/Vehes.htm}
\end{itemize}
\section{Religious and Community Resources}
Joshua Project lists no published materials and no Bible translation: \url{https://joshuaproject.net/people_groups/15779}
\section{Digital and Audio Resources}
\begin{itemize}
\item Global Recordings Network language page: \url{https://globalrecordings.net/en/language/val}. I saw only the listing, not its content.
\item Verbix endangered languages page: \url{https://docs.verbix.com/EndangeredLanguages/Vehes}
\end{itemize}
\section{Curriculum Development Context}
With about 70 speakers and no known written materials, any curriculum would depend on new documentation and community decisions about orthography. Tok Pisin or English would likely be the school languages in practice.
\end{document}